\documentclass[11pt,a4paper]{article}
\usepackage[utf8]{inputenc}
\usepackage[spanish]{babel}
\usepackage[a4paper,left=2.5cm,right=2.5cm,top=2.5cm,bottom=2.5cm]{geometry}
\usepackage{amsmath}
\usepackage{amsfonts}
\usepackage{amssymb}
\usepackage{fancyhdr}
\usepackage{graphicx}
\usepackage{float}
\usepackage{booktabs}
\usepackage{array}
\usepackage{xcolor}
\usepackage{hyperref}
\usepackage{enumitem}
\usepackage{tcolorbox}

\pagestyle{fancy}
\lhead{\small Econométria - Pauta Ejercicios 5}
\rhead{\small Universidad Técnica Federico Santa María}

\definecolor{verdeUSM}{RGB}{0,100,50}

\newtcolorbox{conclusion}{colback=green!5!white,colframe=verdeUSM,title=Conclusión}
\newtcolorbox{nota}{colback=yellow!5!white,colframe=orange!75!black,title=Nota}

\begin{document}

\begin{center}
    {\Large \textbf{Pauta Ejercicios 5: Econométria}}\\[0.3cm]
    Profesor: Pedro Comte Hidalgo \\
    Ayudante: Matias Balboa \\
    \textit{Interpretación aproximada para modelos log-lineales}
\end{center}

\vspace{0.5cm}
\hrule
\vspace{0.5cm}

\section*{Ejercicio 1: Test F de restricciones múltiples}

\textbf{Planteamiento de hipótesis:}
\begin{align*}
H_0 &: \beta_3 = \beta_4 = 0 \quad \text{(las variables no son relevantes)} \\
H_1 &: \text{al menos una es distinta de cero}
\end{align*}

\textbf{Estadístico F basado en \(R^2\):}
\[
F = \frac{(R^2_{UR} - R^2_R)/q}{(1-R^2_{UR})/(n-k-1)}
\]

Donde \(R^2_{UR} = 0.62\), \(R^2_R = 0.55\), \(q = 2\) restricciones, y \(n - k - 1 = 75 - 4 - 1 = 70\) grados de libertad.

\textbf{Cálculo:}
\[
F = \frac{(0.62 - 0.55)/2}{(1-0.62)/70} = \frac{0.035}{0.005429} \approx 6.45
\]

\textbf{Decisión:} Como \(F = 6.45 > F_{0.95,(2,70)} = 3.13\), \textbf{se rechaza \(H_0\)} al 5\% de significancia.

\begin{conclusion}
Las variables \(x_3\) y \(x_4\) son conjuntamente relevantes para explicar \(y\); aportan capacidad explicativa al modelo.
\end{conclusion}

\vspace{0.5cm}
\hrule
\vspace{0.5cm}

\section*{Ejercicio 2: Wagepanel}

\subsection*{(a) Cálculo de \(SSR_{UR}\) y \(SSR_R\)}

Usando la relación \(R^2 = 1 - \frac{SSR}{SST}\), despejamos \(SSR = SST(1-R^2)\):

\[
SSR_{UR} = 200 \cdot (1 - 0.55) = 200 \cdot 0.45 = 90
\]
\[
SSR_R = 200 \cdot (1 - 0.38) = 200 \cdot 0.62 = 124
\]

\subsection*{(b) Test F: \(H_0: \beta_2 = \beta_3 = 0\)}

\[
F = \frac{(R^2_{UR}-R^2_R)/q}{(1-R^2_{UR})/(n-k-1)} = \frac{(0.55-0.38)/2}{(1-0.55)/46} = \frac{0.085}{0.009783} \approx 8.69
\]

\textbf{Decisión:} Como \(F = 8.69 > F_{0.95,2,46} = 3.20\), \textbf{se rechaza \(H_0\)}.

\begin{conclusion}
\textbf{Económicamente:} la experiencia laboral (\(exper\)) y la antigüedad en la empresa (\(tenure\)) son conjuntamente significativas para explicar el logaritmo del salario. La teoría del capital humano se valida: la acumulación de experiencia general y específica es relevante para determinar los salarios.
\end{conclusion}

\subsection*{(c) Test F de significancia global del modelo UR}

\[
F = \frac{R^2_{UR}/k}{(1-R^2_{UR})/(n-k-1)} = \frac{0.55/3}{0.45/46} = \frac{0.1833}{0.009783} \approx 18.74
\]

\textbf{Decisión:} Como \(F = 18.74 > F_{0.95,3,46} = 2.81\), \textbf{se rechaza \(H_0\)}.

\begin{conclusion}
El modelo en su conjunto es estadísticamente significativo; al menos uno de los regresores (educación, experiencia o antigüedad) explica significativamente el logaritmo del salario.
\end{conclusion}

\vspace{0.5cm}
\hrule
\vspace{0.5cm}

\section*{Ejercicio 3: Modelo log-lineal con interacciones y dummies}

\textbf{Modelo estimado:}
\begin{align*}
\ln(wage_i) &= 1.80 - 0.10 \cdot female + 0.070 \cdot educ + 0.010 \cdot (female \times educ) \\
&\quad - 0.050 \cdot parttime - 0.040 \cdot (female \times parttime) + u_i
\end{align*}

\begin{nota}
\textbf{Interpretación aproximada:} como la variable dependiente está en logaritmo, los coeficientes multiplicados por 100 se interpretan \emph{aproximadamente} como cambios porcentuales en el salario. Es decir, si \(\hat{\beta} = 0.07\), entonces un incremento unitario en la variable explicativa se asocia con un aumento aproximado del \(7\%\) en el salario. Esta aproximación es válida para coeficientes pequeños (\(|\hat{\beta}| < 0.15\)).
\end{nota}

\subsection*{(a) Interpretación de \(\hat{\beta}_1\) y \(\hat{\beta}_4\)}

\textbf{\(\hat{\beta}_1 = -0.10\):} Cuando la interacción \(female \times educ\) es cero (para \(educ=0\)), ser mujer se asocia con un salario aproximadamente \textbf{10\% menor} que el de un hombre, manteniendo \(parttime\) constante (cuando \(female \times parttime = 0\)).

\textbf{\(\hat{\beta}_4 = -0.050\):} Cuando la interacción \(female \times parttime\) es cero (para hombres), trabajar a tiempo parcial se asocia con un salario aproximadamente \textbf{5\% menor} que trabajar full-time, manteniendo la educación constante.

\textbf{Sobre la aditividad con dependiente en log:} con \(\ln(wage)\) como dependiente, los efectos se suman \emph{en la escala logarítmica}. Esto significa que, en la escala original del salario, los efectos son \emph{multiplicativos}, no aditivos. Sin embargo, bajo la interpretación aproximada, podemos sumar los efectos porcentuales directamente.

\subsection*{(b) Interpretación de \(\hat{\beta}_3 = 0.010\)}

\(\hat{\beta}_3\) mide la \textbf{diferencia en el retorno de la educación entre mujeres y hombres}.

\begin{itemize}[leftmargin=2em]
    \item Retorno de un año adicional de educación para \textbf{hombres}: \(\hat{\beta}_2 = 0.070 \Rightarrow 7\%\)
    \item Retorno de un año adicional de educación para \textbf{mujeres}: \(\hat{\beta}_2 + \hat{\beta}_3 = 0.080 \Rightarrow 8\%\)
\end{itemize}

Es decir, cada año adicional de educación incrementa el salario aproximadamente \textbf{1 punto porcentual más} para mujeres que para hombres. Las mujeres tienen un retorno marginal a la educación ligeramente mayor.

\subsection*{(c) Efecto de ser mujer para un full-time con \(educ=12\)}

El efecto diferencial total al ser mujer (vs hombre) con \(parttime=0\) y \(educ=12\) es:

\[
\Delta \ln(wage) = \hat{\beta}_1 + \hat{\beta}_3 \cdot educ = -0.10 + 0.010 \cdot 12 = 0.02
\]

\begin{conclusion}
\textbf{Interpretación aproximada:} una mujer full-time con 12 años de educación gana aproximadamente \textbf{2\% más} que un hombre comparable. Aunque la brecha inicial es negativa (\(-10\%\)), el mayor retorno a la educación de las mujeres (\(+1\%\) por año) la compensa a partir de cierto nivel educativo.
\end{conclusion}

\subsection*{(d) Efecto de ser part-time para hombres y mujeres}

\textbf{Para hombres} (con \(female=0\)):
\[
\Delta \ln(wage) = \hat{\beta}_4 = -0.050 \Rightarrow \text{aproximadamente } -5\%
\]
Un hombre part-time gana aproximadamente \textbf{5\% menos} que un hombre full-time.

\textbf{Para mujeres} (con \(female=1\)):
\[
\Delta \ln(wage) = \hat{\beta}_4 + \hat{\beta}_5 = -0.050 - 0.040 = -0.090 \Rightarrow \text{aproximadamente } -9\%
\]
Una mujer part-time gana aproximadamente \textbf{9\% menos} que una mujer full-time.

\begin{conclusion}
La penalización salarial por trabajar part-time es mayor para mujeres (\(-9\%\)) que para hombres (\(-5\%\)).
\end{conclusion}

\subsection*{(e) Ecuaciones de predicción por subgrupo}

\textbf{Hombre full-time} (\(female=0\), \(parttime=0\)):
\[
\widehat{\ln(wage)} = 1.80 + 0.070 \cdot educ
\]

\textbf{Hombre part-time} (\(female=0\), \(parttime=1\)):
\[
\widehat{\ln(wage)} = 1.75 + 0.070 \cdot educ
\]

\textbf{Mujer full-time} (\(female=1\), \(parttime=0\)):
\[
\widehat{\ln(wage)} = 1.70 + 0.080 \cdot educ
\]

\textbf{Mujer part-time} (\(female=1\), \(parttime=1\)):
\[
\widehat{\ln(wage)} = 1.61 + 0.080 \cdot educ
\]

\vspace{0.3cm}

\textbf{Resumen de cambios:}

\begin{center}
\begin{tabular}{lcc}
\toprule
\textbf{Subgrupo} & \textbf{Intercepto} & \textbf{Pendiente en educ} \\
\midrule
Hombre full-time & 1.80 & 0.070 (\(7\%\)) \\
Hombre part-time & 1.75 & 0.070 (\(7\%\)) \\
Mujer full-time  & 1.70 & 0.080 (\(8\%\)) \\
Mujer part-time  & 1.61 & 0.080 (\(8\%\)) \\
\bottomrule
\end{tabular}
\end{center}

\textbf{Observaciones clave:}
\begin{itemize}[leftmargin=2em]
    \item \textbf{Cambio de pendiente en educación:} solo depende del género. Las mujeres tienen pendiente \(0.080\) y los hombres \(0.070\) (diferencia de 1 pp).
    \item \textbf{Cambio de nivel por part-time:} \(-0.050\) para hombres y \(-0.090\) para mujeres (es decir, \(-0.050 + (-0.040)\) por la interacción).
\end{itemize}

\end{document}